\documentclass[11pt, letterpaper]{article}
\usepackage[T1]{fontenc}
\usepackage{amsmath, amssymb, amsthm}
\usepackage[margin=1.12in]{geometry}
\usepackage{lmodern}
\usepackage[hidelinks]{hyperref}
\hypersetup{
    pdftitle={Hereditary atomicity and hereditary ACCP in semilocal domains},
    pdfauthor={Astra 6.2},
    pdfsubject={AI-authored research manuscript on factorization in subrings of semilocal domains}
}
\newtheorem{maintheorem}{Theorem}
\renewcommand{\themaintheorem}{\Alph{maintheorem}}
\newtheorem{theorem}{Theorem}[section]
\newtheorem{lemma}[theorem]{Lemma}
\newtheorem{proposition}[theorem]{Proposition}
\newtheorem{corollary}[theorem]{Corollary}
\theoremstyle{definition}
\newtheorem{definition}[theorem]{Definition}
\theoremstyle{remark}
\newtheorem{remark}[theorem]{Remark}
\DeclareMathOperator{\Frac}{Frac}
\raggedbottom

\title{Hereditary atomicity and hereditary ACCP in semilocal domains}
\author{AI-authored manuscript}
\date{November 17, 2026}

\begin{document}
\maketitle

\begin{abstract}
    For a semilocal integral domain $D$, every unital subring of $D$ is atomic if and only if every such subring satisfies the ascending chain condition on principal ideals (ACCP).
    Under these conditions $D_M$ also satisfies ACCP for every maximal ideal $M$ of $D$.
    The local construction starts with a local domain that contains a discrete valuation ring and has fraction field algebraic over that of the base.
    From a strict chain of principal ideals we construct a nonatomic subring containing the base.
    In the semilocal case we lift generators of a chain at one maximal ideal to elements that are units at all the others, so the obstruction to factorization survives inside $D$.
    For a nonfield semilocal hereditarily atomic domain we then bound the number of factors of a fixed element that can be nonunits of a subring while being units of $D$.
\end{abstract}

\section{Introduction}
\label{sec:intro}

Suppose a nonzero nonunit $a$ of a domain $R$ has no factorization into irreducibles.
Then $a$ splits as a product of two nonunits, at least one of which again has no such factorization.
Repeating this choice produces a strictly ascending chain of principal ideals beginning with $aR$.
The converse fails even for local domains.
Grams constructed a local atomic domain with an infinite strictly ascending chain of principal ideals \cite[Theorem~1.3]{Grams}.
For such a domain the question becomes whether passage to a subring can destroy all the irreducible factorizations of a single element.

All rings below are commutative with identity, and subrings contain the same identity.
An \emph{atom} is a nonzero nonunit that cannot be expressed as a product of two nonunits.
A domain is \emph{atomic} if every nonzero nonunit is a finite product of atoms, and \emph{hereditarily atomic} if every subring is atomic.
It satisfies the \emph{ascending chain condition on principal ideals} (ACCP) if every ascending sequence of principal ideals stabilizes.
A domain is \emph{hereditarily ACCP} if every subring satisfies ACCP.
A domain is \emph{semilocal} if it has finitely many maximal ideals.
Coykendall, Gotti and Hasenauer formulated the conjecture that hereditary atomicity implies hereditary ACCP and showed that this amounts to asking whether every domain that fails ACCP contains a nonatomic subring \cite[Conjecture~6.1 and Proposition~6.2]{CGH}.
For fields they proved that in every subring of a hereditarily atomic field, the lengths of the factorizations of a fixed element into atoms are bounded \cite[Corollary~4.5]{CGH}.
We prove that hereditary atomicity and hereditary ACCP are equivalent for semilocal domains.
We begin with the local case, where a single chain of principal ideals provides the construction.

\begin{maintheorem}
    \label{thm:local}
    Every local domain that fails ACCP contains a nonatomic subring.
    Equivalently, every hereditarily atomic local domain satisfies ACCP.
\end{maintheorem}

The construction preserves a specified coefficient ring in the following setting.

\begin{maintheorem}
    \label{thm:dvr}
    Let $\mathcal{O}$ be a discrete valuation ring (DVR) contained in a local domain $R$.
    Suppose $\Frac(R)$ is algebraic over $\Frac(\mathcal{O})$.
    If $R$ fails ACCP, there is a nonatomic subring $S$ with $\mathcal{O} \subseteq S \subseteq R$.
\end{maintheorem}

For a semilocal domain the conclusion extends to subrings with arbitrarily many maximal ideals.

\begin{maintheorem}
    \label{thm:hereditary}
    A semilocal domain $D$ is hereditarily atomic if and only if it is hereditarily ACCP.
    Under these equivalent conditions $D_M$ satisfies ACCP for every maximal ideal $M$ of $D$.
\end{maintheorem}

Section~\ref{sec:arbitrary-local} deduces Theorem~\ref{thm:local} from Theorem~\ref{thm:dvr}.
The two possibilities in the deduction explain where algebraic relations enter the proof.
A transcendental element can serve as a polynomial variable, and taking its constant term produces a prime ideal containing no atoms.
In the algebraic case, start with a strict chain $x_{0}R \subsetneq x_{1}R \subsetneq \cdots$.
Group consecutive quotients $x_{i - 1}/x_i$ into products $y_j$ and scale the generators at the block endpoints to obtain elements $z_j$ satisfying $z_{j - 1} = y_j z_j$.
We construct an auxiliary subring $R_{*}$ and an ideal
\[
    J_{*} = \bigcup_{j \ge 1} z_j R_{*}.
\]
Adjoining $J_{*}$ to the ring generated over $\mathcal{O}$ by a tail of the $y_j$ gives our subring $S$.
Each nonzero element of $J_{*}$ is divisible in $S$ by a later $y_j$, with the quotient still in $J_{*}$.
The remaining task is to show that every factorization of $z_1$ has a factor in $J_{*}$.

One cannot deduce this from primality.
In fact a uniformizer of $\mathcal{O}$ lies outside $J_{*}$ while one of its powers lies inside, as we show in Remark~\ref{rem:nonprime}.
We therefore study the particular element $z_1$ in a Laurent polynomial ring with independent variables $Y_j$ representing the $y_j$.
Any proposed factorization uses only finitely many blocks.
We choose successive blocks of sufficiently high algebraic degree that the relations introduced by $Y_j \mapsto y_j$ contribute only a controlled error to that finite calculation.
Theorem~\ref{thm:F} excludes a factorization even in the presence of this error.

For a semilocal domain $D$, Theorem~\ref{thm:hereditary} gives a nonatomic subring of $D$ whenever $D_M$ fails ACCP at even one maximal ideal $M$.
When $D_M$ satisfies the base hypotheses of Proposition~\ref{prop:semilocal-witness}, we construct the subring by lifting the generators of the local chain to units at all the other maximal ideals.
This choice ensures that the divisions used to split factors still take place inside $D$.

The passage to an arbitrary subring $S \subseteq D$ in Theorem~\ref{thm:hereditary} involves a different obstruction.
Even after the induced chain in $D$ stabilizes, a quotient that is a unit of $D$ can still be a nonunit of $S$.
We detect those quotients using valuation rings of $\Frac(D)$ that contain $S$ and omit some element of $D$.
For nonfield semilocal hereditarily atomic $D$, Section~\ref{sec:places} proves that these valuation rings are discrete and that only finitely many can assign positive value to the first generator of the chain.
Their normalized values on that generator bound the total number of quotients that become units.
Section~\ref{sec:hereditary} uses this bound to prove hereditary ACCP and deduces the field case from the nonfield case.

\section{Passing from a DVR base to an arbitrary local domain}
\label{sec:arbitrary-local}

A suitable DVR inside the given local domain will put us in the setting of Theorem~\ref{thm:dvr} whenever the fraction field extension is algebraic.
The following observation treats the transcendental case.

\begin{lemma}
    \label{lem:transcendental-witness}
    Let $A \subseteq R$ be domains and suppose $q \in A$ is a nonunit of $A$ with $q^{-1} \in R$.
    For each $x \in R$ transcendental over $\Frac(A)$, the subring
    \[
        S = A + xA[q^{-1}][x]
    \]
    of $R$ is nonatomic.
\end{lemma}

\begin{proof}
    Put $B = A[q^{-1}]$.
    Since $x$ is transcendental, taking constant terms gives a surjection $S \to A$ whose kernel $I = xB[x]$ is prime.
    The image of $q$ is a nonunit, so $q$ is a nonunit in $S$ as well.
    Every nonzero $f \in I$ factors as $q(f/q)$ with $f/q \in I$.
    Both factors are nonunits of $S$, so $I$ contains no atom.
    Any factorization of $x$ into atoms would put an atom in $I$ by primality.
\end{proof}

\begin{proof}[Proof of Theorem~\ref{thm:local}]
    Let $(R, \mathfrak{m})$ be a local domain failing ACCP.
    In particular $\mathfrak{m} \ne 0$.
    We exhibit a principal ideal domain $A \subseteq R$ and a prime element $\pi$ of $A$ such that
    \[
        A \cap \mathfrak{m} = \pi A,
        \qquad 1 + \pi \notin A^{\times}.
    \]
    If $R$ has characteristic zero and its residue field has characteristic $p > 0$, take $A = \mathbb{Z}$ and $\pi = p$.
    The map $\mathbb{Z} \to R/\mathfrak{m}$ has kernel $p\mathbb{Z}$, and $1 + p > 1$.

    In equal characteristic, $R$ contains its prime field $k$.
    Indeed, in residue characteristic zero every nonzero integer lies outside $\mathfrak{m}$ and hence is invertible in $R$.
    Choose nonzero $\pi \in \mathfrak{m}$.
    If $\pi$ were algebraic over $k$, then $k[\pi]$ would be a field, making $\pi$ invertible in $R$.
    Thus $A = k[\pi]$ is a polynomial ring in one variable, and $1 + \pi$ is a nonunit of $A$.
    Reduction modulo $\mathfrak{m}$ sends a polynomial in $\pi$ to its constant coefficient in the embedded field $k$, so its kernel is $\pi A$.

    In either case $\mathcal{O} = A_{(\pi)}$ lies in $R$, since $R$ inverts every element of $A \setminus \pi A$.
    It is a DVR with fraction field $F = \Frac(A)$.
    If $\Frac(R)$ is algebraic over $F$, Theorem~\ref{thm:dvr} supplies a nonatomic subring of $R$ containing $\mathcal{O}$.
    Otherwise some $x \in R$ is transcendental over $F$, since the elements algebraic over $F$ in $\Frac(R)$ form a field.
    The element $q = 1 + \pi$ is a unit of $R$ because it is congruent to $1$ modulo $\mathfrak{m}$.
    Applying Lemma~\ref{lem:transcendental-witness} gives a nonatomic subring of $R$ in this case as well.
\end{proof}

\section{Normalizing a chain over a discrete valuation ring}
\label{sec:chain}

Fix a discrete valuation ring $\mathcal{O}$ with uniformizer $\varpi$ and fraction field $k_{0}$.
To prove Theorem~\ref{thm:dvr} we first normalize a strict chain so that its initial element is a power of $\varpi$.
The following data specify everything about the ambient ring used in the construction.
Valuations are additive, with $v(0) = +\infty$.

\begin{definition}
    \label{hyp:M}
    A \emph{normalized chain} consists of the following data.
    \begin{enumerate}
        \item A field $K$ algebraic over $k_{0}$ and a valuation ring $V$ of $K$ extending $\mathcal{O}$, with real valuation $v$ normalized by $v(\varpi) = 1$.
        \item An integer $\kappa \ge 1$ and elements $c_{i} \in K^{\times}$ for $i \ge 0$, with $c_{0} = \varpi^{\kappa}$, $v(c_{i}) \ge 1$, and successive quotients $t_{i} = c_{i - 1}/c_{i}$ of positive value for $i \ge 1$.
        \item A subring $A$ of $K$ containing $\mathcal{O}$, all $c_{i}$ and all $t_{i}$, and integers $n_{i} \ge 1$ with $t_{i}^{n_{i}} \in \varpi^{\kappa} A$ for $i \ge 1$.
    \end{enumerate}
\end{definition}

\begin{lemma}
    \label{lem:normalize}
    Under the hypotheses of Theorem~\ref{thm:dvr}, every strictly increasing chain of nonzero principal ideals $x_{i}R$ gives a normalized chain with $A = R$ and $c_{i} = sx_{i}$ for a single nonzero $s \in R$.
\end{lemma}

\begin{proof}
    First, every nonzero ideal $I$ of $R$ contains a power of $\varpi$.
    Indeed, choose nonzero $a \in I$ and clear denominators in its minimal polynomial over $k_{0}$.
    The nonzero constant term belongs to $aR \cap \mathcal{O}$ and is a unit of $\mathcal{O}$ times a power of $\varpi$.
    In particular every nonzero prime $P$ of $R$ contains $\varpi$.
    For any $b \in R$ take a polynomial relation over $\mathcal{O}$ for $b$ and divide its coefficients by their largest common power of $\varpi$.
    Its image modulo $P$ is a nonzero polynomial over $\mathcal{O}/\varpi\mathcal{O}$ annihilating the image of $b$.
    Thus $R/P$ is a domain algebraic over a field, hence a field.
    It follows that the maximal ideal $\mathfrak{m}$ of $R$ is its only nonzero prime ideal.

    The strict inclusions give $t_{i} = x_{i - 1}/x_{i} \in \mathfrak{m}$ and $x_{0} \in \mathfrak{m}$.
    Apply the ideal assertion to $x_{0}R$ to write $wx_{0} = \varpi^{r}$ with $w \in R$ and $r \ge 1$.
    Set $s = \varpi w$, $c_{i} = sx_{i}$, and $\kappa = r + 1$.
    Since $\sqrt{\varpi^{\kappa}R} = \mathfrak{m}$, each $t_{i}$ has a positive power in $\varpi^{\kappa}R$.

    Take a valuation ring $V$ of $K = \Frac(R)$ dominating $R$ \cite{Stacks}.
    Its maximal ideal contracts to $\mathfrak{m}$, so $\varpi$ has positive value and $V \cap k_{0} = \mathcal{O}$.
    The value group embeds in $\mathbb{Q}$ after assigning $\varpi$ value one.
    Indeed, in a polynomial relation for any $b \in K^{\times}$ over $k_{0}$, two nonzero terms must have equal minimum value.
    Subtracting their values shows that a positive multiple of the value of $b$ is an integral multiple of the value of $\varpi$.
    In this normalization $v(t_{i}) > 0$ and $v(c_{i}) = 1 + v(wx_{i}) \ge 1$, as required.
\end{proof}

Fix a normalized chain for the rest of the proof.
We shall construct the witness inside $A \cap V$.
For a vector $\beta$ write $\beta_{\le i}$ and $\beta_{>i}$ for the coordinates with indices at most $i$ and greater than $i$, respectively.
Inequalities between vectors are coordinatewise.

\section{Coordinates for products of chain ratios}
\label{sec:windows}

Write $\widehat{K}$ for the completion of $K$ at $v$ and $\widehat{k}_{0}$ for the closure of $k_{0}$ in $\widehat{K}$.
The latter is complete and discretely valued.
Indeed, an approximation to a nonzero element by an element of $k_{0}$ has the same value once the error has larger value.
We denote its valuation ring by $\widehat{\mathcal{O}}$.
Every finite extension of $\widehat{k}_{0}$ in $\widehat{K}$ is complete, and its valuation extends uniquely to its algebraic closure \cite[Proposition~1.2.2 and the discussion in \S4.1]{MR2183496}.

\begin{lemma}
    \label{lem:L}
    Let $F$ be a finite extension of $\widehat{k}_{0}$ in $\widehat{K}$, and let $y \in \widehat{K}^{\times}$ be algebraic of degree $D$ over $F$.
    Fix any interval $I$ of $D$ consecutive integers.
    In the basis $(y^{j})_{j \in I}$ of $F(y)$ over $F$, each term of the expansion of $y^{n}$ has value at least $n v(y)$, for every integer $n$.
\end{lemma}

\begin{proof}
    Put $u = v(y)$ and write the minimal polynomial as $\sum_{j = 0}^{D} a_{j} X^{j}$, with $a_{D} = 1$.
    Uniqueness of the extended valuation gives every conjugate of $y$ value $u$.
    Taking elementary symmetric functions of the roots with their multiplicities gives $v(a_{j}) \ge (D - j)u$ and $v(a_{0}) = Du$.
    Start with the usual power basis.
    To multiply a basis expansion by $y$, one need only reduce $y^{D}$ using the minimal polynomial, and every replacement term has value at least $Du$.
    To multiply by $y^{-1}$, use $y^{-1} = -\sum_{j = 1}^{D}(a_{j}/a_{0})y^{j - 1}$, whose terms have value at least $-u$.
    Thus successive multiplications starting at $1$ give the bound for every $n$, since collecting equal powers can only raise their value.
    If $h = \min I$, multiplying the expansion of $y^{n - h}$ by $y^{h}$ proves the assertion in the translated basis.
\end{proof}

Consider a tower $F_{0} = \widehat{k}_{0}$, $F_{i} = F_{i - 1}(y_{i})$ generated by nonzero elements of $K$.
Each step is finite since $K/k_{0}$ is algebraic.
Put $D_{i} = [F_{i} : F_{i - 1}]$ and $u_{i} = v(y_{i})$.
The transfer to Laurent polynomials will keep a monomial's final coordinates fixed while expanding its initial coordinates in a basis.
For this reason the interval for an exponent depends on the poles in the coordinates to its right.

\begin{definition}
    \label{def:WN}
    Put $L_{i} = \lfloor D_{i}/2 \rfloor$ and $I_{i} = [-L_{i}, D_{i} - L_{i} - 1] \cap \mathbb{Z}$.
    For $\beta \in \mathbb{Z}^{N}$ let $\delta_{i}(\beta) = \max\bigl(\{0\} \cup \{-\beta_{j} : i < j \le N\}\bigr)$.
    A vector $\beta$ fails the \emph{window condition} at $i$ if $\beta_{i} + \delta_{i}(\beta) \notin I_{i}$.
    Write $B_{N}$ for the set of vectors that meet all the interval conditions and $B_{0}$ for the singleton consisting of the empty vector.
\end{definition}

\begin{lemma}
    \label{lem:W}
    Fix the last $N - l$ coordinates of $\beta$, and suppose they meet their interval conditions.
    Put $\delta = \delta_{l}(\beta)$.
    Then $\beta \in B_{N}$ precisely when $\beta_{\le l} + \delta\mathbf{1} \in B_{l}$, where $\mathbf{1}$ denotes the vector of ones.
\end{lemma}

\begin{proof}
    For $\gamma = \beta_{\le l} + \delta\mathbf{1}$ and $i \le l$, the largest pole order to the right of $i$ is $\delta_{i}(\beta) = \delta + \delta_{i}(\gamma)$.
    This follows by taking the maximum first over $j > l$ and then over $i < j \le l$.
    Hence $\gamma_{i} + \delta_{i}(\gamma) = \beta_{i} + \delta_{i}(\beta)$ at every remaining coordinate.
\end{proof}

\begin{lemma}
    \label{lem:WB}
    The monomials $y^{\beta} = \prod_{i = 1}^{N} y_{i}^{\beta_{i}}$ for $\beta \in B_{N}$ form a basis of $F_{N}$ over $\widehat{k}_{0}$.
\end{lemma}

\begin{proof}
    At $N = 0$ the basis is $\{1\}$.
    For $N \ge 1$, decompose $F_{N}$ into the summands $F_{N - 1} y_{N}^{j}$ for $j \in I_{N}$ by Lemma~\ref{lem:L}.
    In the summand for $j$, put $d = \max(0, -j)$ and multiply the preceding basis by $(y_{1}\cdots y_{N - 1})^{-d} y_{N}^{j}$.
    This gives a basis of the summand, and Lemma~\ref{lem:W} identifies the resulting exponent vectors as exactly the elements of $B_{N}$ ending in $j$.
\end{proof}

\begin{lemma}
    \label{lem:WL}
    There are nondecreasing constants $\Omega_{1} = 0, \Omega_{2}, \dots$ such that every term of the $B_{l}$-expansion of $x \in F_{l}$ has value at least $v(x) - \Omega_{l + 1}$.
    The constant $\Omega_{l + 1}$ depends only on the generators chosen through stage $l$.
    For a single monomial $x = c y^{\alpha}$ with $c \in \widehat{k}_{0}^{\times}$, $\alpha \in \mathbb{Z}^{l}$ and $l \ge 1$, the lower bound improves to $v(x) - \Omega_{l}$.
\end{lemma}

\begin{proof}
    For $|a| = \exp(-v(a))$, the maximum of the absolute values of the basis terms is a norm on $F_{l}$.
    Equivalence of norms over the complete field $\widehat{k}_{0}$ bounds this maximum by a constant times $|x|$ \cite[Proposition~1.2.2]{MR2183496}.
    Take $\Omega_{l + 1}$ to be the least nonnegative logarithmic bound.
    Equivalently, take the supremum of $v(x) - v(t)$ over all $x \ne 0$ and all nonzero terms $t$ in the basis expansion of $x$.
    It is zero for $l = 0$, and it can only increase with $l$ because adjoining the exponent $0$ embeds each basis into the next.

    Now write a monomial $x$ as $a y_{l}^{n}$ with $a \in F_{l - 1}^{\times}$.
    Expansion in the last variable gives $x = \sum_{j \in I_{l}} g_{j} y_{l}^{j}$ with $v(g_{j} y_{l}^{j}) \ge v(x)$ by Lemma~\ref{lem:L}.
    For each $j$, set $b = (y_{1} \cdots y_{l - 1})^{\max(0, -j)}$.
    Expand $b g_{j}$ in the preceding basis and multiply the expansion by $y_{l}^{j}/b$.
    Its terms have value at least $v(g_{j} y_{l}^{j}) - \Omega_{l}$.
    By the summand decomposition in Lemma~\ref{lem:WB}, these are precisely the required basis terms.
\end{proof}

\section{Selecting the blocks}
\label{sec:blocks}

Return to the normalized chain of Definition~\ref{hyp:M}.
At the $j$th stage, re-expanding a single monomial can lower the value of a term by as much as $\Omega_{j}$.
We can offset this loss by selecting a sufficiently long product of successive ratios, because its field degree can be made arbitrarily large.

\begin{lemma}
    \label{lem:degreegrowth}
    Fix $a \ge 0$ and a finite extension $F$ of $\widehat{k}_{0}$ in $\widehat{K}$.
    The degrees $[F(c_{a}/c_{b}) : F]$ for $b > a$ are unbounded.
\end{lemma}

\begin{proof}
    As $b$ increases, the values $v(c_{a}/c_{b})$ increase strictly and stay below $v(c_{a})$, by the chain hypotheses.
    If the indicated degrees were bounded, the degrees over $\widehat{k}_{0}$ of the elements $c_{a}/c_{b}$ would all be at most some integer $M$.
    An element of degree $d$ over $\widehat{k}_{0}$ has value in $d^{-1}\mathbb{Z}$, since the constant coefficient of its monic minimal polynomial has value $d$ times the value of the element, as in Lemma~\ref{lem:L}.
    Consequently the strictly increasing sequence would be bounded in $(M!)^{-1}\mathbb{Z}$, which is impossible.
\end{proof}

\begin{proposition}
    \label{prop:B}
    There are integers $0 = a_{0} < a_{1} < \cdots$ with the following properties.
    Set $z_{0} = \varpi^{\kappa}$, $s_{0} = \kappa$, $T_{0} = 0$, and for $j \ge 1$ put
    \[
        z_{j} = c_{a_{j}}, \qquad y_{j} = z_{j - 1}/z_{j}, \qquad
        s_{j} = v(z_{j}), \qquad u_{j} = v(y_{j}).
    \]
    Use these $y_{j}$ to form the tower above, and set $m_{j} = n_{a_{j - 1} + 1}$, using the exponents from the chain hypotheses.
    Then $y_{j}, z_{j} \in A$, and the following hold for every $j \ge 1$.
    \begin{align}
        s_{j} &= \kappa - \sum_{i = 1}^{j} u_{i} \ge 1,
        \qquad 0 < u_{j} < \kappa, \tag{H1}\label{H1}\\
        T_{j} &= T_{j - 1} + 2\kappa m_{j} \ge T_{j - 1} + 2\kappa, \tag{H2}\label{H2}\\
        D_{j}u_{j} &\ge 3\kappa(2T_{j} + \Omega_{j}). \tag{H3}\label{H3}
    \end{align}
    Write $\tau_{j} = T_{j} + s_{j}$ for $j \ge 0$.
\end{proposition}

\begin{proof}
    Once $a = a_{j - 1}$ is chosen, the numbers $m_{j}$, $T_{j}$ of \eqref{H2}, and $\Omega_{j}$ of Lemma~\ref{lem:WL} are fixed.
    By Lemma~\ref{lem:degreegrowth} choose an integer $a_{j} > a$ such that $[F_{j - 1}(c_{a}/c_{a_{j}}) : F_{j - 1}]\,v(t_{a + 1}) \ge 3\kappa(2T_{j} + \Omega_{j})$.
    The new ratio $y_{j} = \prod_{a < i \le a_{j}} t_{i}$ lies in $A$ and has value at least $v(t_{a + 1}) > 0$, proving \eqref{H3}.
    Since $z_{j}$ is a term of the original chain, it belongs to $A$ and has value at least one.
    Telescoping gives \eqref{H1}, and $m_{j} \ge 1$ gives the inequality in \eqref{H2}.
\end{proof}

\begin{lemma}
    \label{lem:D}
    For nonnegative integers $e, q_{1}, \dots, q_{l}$ with $l \ge 1$, the element $\varpi^{e - \kappa}\prod_{i = 1}^{l} y_{i}^{q_{i}}$ belongs to $A$ whenever its value is at least $T_{l}$.
\end{lemma}

\begin{proof}
    The first ratio of block $i$ has $m_{i}$th power in $\varpi^{\kappa}A$.
    All remaining ratios belong to $A$, so $y_{i}^{m_{i}}/\varpi^{\kappa} \in A$.
    If $e \ge \kappa$ the given element is already a product in $A$.
    For $e < \kappa$, if all $q_{i} < m_{i}$, the element has value less than $\kappa \sum_{i = 1}^{l} m_{i} = T_{l}/2$ by \eqref{H1} and \eqref{H2}.
    The assumed lower bound therefore gives some $q_{i} \ge m_{i}$, and the factor $y_{i}^{m_{i}}$ clears the denominator.
\end{proof}

\section{A finite Laurent polynomial model}
\label{sec:model}

Replace a finite string of blocks by independent variables of the same values.
The relations $Z_{k - 1} = Y_{k}Z_{k}$ will then hold identically.
The following rings record which Laurent monomials may occur on each side of a factorization.
We say membership in a set of Laurent polynomials is \emph{termwise} when a polynomial belongs to the set precisely when all its nonzero terms belong.

\begin{definition}
    \label{def:GS}
    Work over the completed discretely valued field $\widehat{k}_{0}$ with valuation ring $\widehat{\mathcal{O}}$ and $v(\varpi) = 1$.
    Fix positive integers $N, \kappa$ and positive real weights $u_{1}, \ldots, u_{N}$ with $s_{k} = \kappa - \sum_{i = 1}^{k} u_{i} > 0$ for $0 \le k \le N$.
    Take $T_{0} = 0$, $T_{k} \ge 2\kappa$ for $k \ge 1$, and $\tau_{k} = T_{k} + s_{k}$.
    For a set $X \subseteq \{1, \ldots, N\}$, write $u_{X} = \sum_{i \in X} u_{i}$, $Y_{X} = \prod_{i \in X} Y_{i}$ and $X_{\le k} = X \cap \{1, \ldots, k\}$.
    Fix $\Lambda \subseteq \{1, \ldots, N\}$ with $2u_{\Lambda \setminus \{1\}} \le s_{1}$.
    In the Laurent polynomial ring $\mathcal{L} = \widehat{k}_{0}[Y_{1}^{\pm 1}, \ldots, Y_{N}^{\pm 1}]$, put
    \[
        Z_{k} = \varpi^{\kappa}\prod_{i = 1}^{k}Y_{i}^{-1}
        \qquad (0 \le k \le N).
    \]
    For a monomial $t = cY^{\beta}$ with $c \ne 0$, set $e(t) = v(c)$ and $v(t) = e(t) + \beta \cdot u$.
    For a nonzero Laurent polynomial $f$, define $v(f)$ to be the least value of its terms after like powers are collected, and put $v(0) = +\infty$.
    When exactly one term attains the minimum, write $\operatorname{in}(f)$ for that term.
\end{definition}

\begin{lemma}
    \label{lem:V}
    The function $v$ is a valuation on $\mathcal{L}$.
    If a product has a unique term of least value, so does each factor and each subproduct.
    The exponents and the coefficient values of the initial terms add under these products.
\end{lemma}

\begin{proof}
    Among terms of least value, take the first exponent in lexicographic order.
    If $f$ and $g$ have first exponents $a$ and $b$, only this pair of least-value terms can contribute at $a + b$.
    All other contributions at that exponent have higher coefficient valuation, so the product has value $v(f) + v(g)$ and first exponent $a + b$.
    The valuation inequality for sums follows coefficient by coefficient.
    Reversing the order gives additivity of the last exponent as well.
    The difference between the last and first exponents is lexicographically nonnegative for every factor.
    If the differences sum to zero, they all vanish.
    This proves uniqueness in each factor and subproduct, and subtracting the weight of the exponent from the total value gives additivity of the coefficient values.
\end{proof}

\begin{definition}
    \label{def:GM}
    For $t = cY^{\beta}$ let $\delta_{k}(t) = \max(\{0\} \cup \{-\beta_{i} : i > k\})$.
    Put
    \[
        P_{k}(t) = e(t) + \sum_{i = 1}^{k} \beta_{i}u_{i},
        \qquad
        \varphi_{k}(t) = P_{k}(t) - \delta_{k}(t)s_{k}.
    \]
    Write $\mathcal{H}_{k}$ for the Laurent monomials in the first $k$ variables of value at least $T_{k}$, for $1 \le k \le N$.
    Define
    \[
        \begin{aligned}
            R_{0} &= \widehat{\mathcal{O}}[Y_{1}, \ldots, Y_{N}, Z_{N},
                        \mathcal{H}_{1} \cup \cdots \cup \mathcal{H}_{N}], \\
            J_{0} &= Z_{N}R_{0}, \qquad
            S_{0} = \widehat{\mathcal{O}}[Y_{i} : i \in \Lambda] + J_{0}.
        \end{aligned}
    \]
    For a polynomial monomial with $e(t) \ge 0$ and $\beta \ge 0$, say that it is \emph{dead at level $k$} if $P_{k}(t) \ge \tau_{k}$.
    At level zero this means $e(t) \ge \kappa$.
\end{definition}

\begin{lemma}
    \label{lem:MC}
    Suppose a class of monomials is closed under replacing a coefficient by one of equal or greater value.
    Its additive span consists precisely of the polynomials whose terms all belong to the class.
    Taking the union of two such classes adds their spans, and translation by a fixed monomial multiplies the span by that monomial.
\end{lemma}

\begin{proof}
    At each exponent the coefficient in a sum, if nonzero, has valuation at least the smallest among the contributing coefficients.
    Conversely, a polynomial is the sum of its terms.
    Translation by a monomial preserves the distinctness of the exponents.
\end{proof}

\begin{lemma}
    \label{lem:G1}
    For Laurent monomials $t = cY^{\beta}$ and $t'$ the following hold.
    \begin{enumerate}
        \item[\textup{(a)}] The values $\varphi_{k}(t)$ increase weakly to $\varphi_{N}(t) = v(t)$, and $\varphi_{k}(tt') \ge \varphi_{k}(t) + \varphi_{k}(t')$.
        \item[\textup{(b)}] Membership in $R_{0}$ is termwise.
        A monomial $t$ belongs to $R_{0}$ precisely when $\varphi_{k}(t) \ge T_{k}$ at some level $0 \le k \le N$.
        At any such level, with $d = \delta_{k}(t)$, it factors as
        \[
            t = h Z_{N}^{d}\prod_{i > k}Y_{i}^{\beta_{i} + d},
        \]
        where the displayed exponents $\beta_{i} + d$ are nonnegative and $h \in \mathcal{H}_{k}$ for $k > 0$, or $h \in \widehat{\mathcal{O}}$ for $k = 0$.
    \end{enumerate}
    In particular $v \ge 0$ on $R_{0}$, and $Z_{k} \in R_{0}$ for all $k$.
\end{lemma}

\begin{proof}
    Write $d_{k} = \delta_{k}(t)$.
    Since $d_{k - 1} = \max\{d_{k}, -\beta_{k}\}$, the increment of $\varphi$ is
    \[
        \varphi_{k}(t) - \varphi_{k - 1}(t)
        = (\beta_{k} + d_{k - 1})u_{k} + (d_{k - 1} - d_{k})s_{k} \ge 0.
    \]
    Now $d_{N} = 0$, and the pole orders are subadditive under products, whereas the partial values are additive.
    This gives~(a).

    For (b), the monomials passing the stated test are closed under raising the coefficient value.
    If two monomials pass at different levels, use the larger level in~(a).
    There one factor already has the required value and the other contributes a nonnegative value.
    Their additive span is therefore a ring by Lemma~\ref{lem:MC}.
    It contains the generators of $R_{0}$, since constants in $\widehat{\mathcal{O}}$, the $Y_{i}$ and $Z_{N}$ pass at zero, and the monomials of $\mathcal{H}_{k}$ pass at $k$.
    Conversely, if $t$ passes at $k$, set $h = c\varpi^{-\kappa d}\prod_{i \le k}Y_{i}^{\beta_{i} + d}$.
    Then $h$ has value $\varphi_{k}(t)$ and gives the asserted factorization in $R_{0}$.
    This proves the termwise criterion.
    Every passing monomial has nonnegative value by (a), and $Z_{k} = Z_{N}\prod_{i > k}Y_{i}$.
\end{proof}

\begin{lemma}
    \label{lem:G2}
    The following tests hold in the finite model.
    \begin{enumerate}
        \item[\textup{(a)}] Membership in $J_{0}$ is termwise, and $t \in J_{0}$ precisely when $t/Z_{N} \in R_{0}$.
        Every nonzero element of $J_{0}$ has value at least $s_{N}$, and $Z_{k}R_{0} \subseteq J_{0}$.
        \item[\textup{(b)}] For a monomial $t$ with nonnegative exponents, $t \in J_{0}$ precisely when $P_{k}(t) \ge \tau_{k}$ for some $0 \le k \le N$.
        \item[\textup{(c)}] Membership in $S_{0}$ is termwise.
        Each term of an element of $S_{0}$ lies in $J_{0}$ or in $\widehat{\mathcal{O}}[Y_{i} : i \in \Lambda]$.
        A term outside $J_{0}$ is dead at no level.
        \item[\textup{(d)}] If a monomial $t \in J_{0}$ has $v(t) < 2\kappa$, then $e(t) \ge \kappa$.
        In case of equality $t = cZ_{N}Y^{\beta}$ with $v(c) = 0$ and $\beta \ge 0$.
    \end{enumerate}
\end{lemma}

\begin{proof}
    Multiplication by $Z_{N}$ shifts all exponents and adds $s_{N}$ to every value.
    Part~(a) follows from Lemma~\ref{lem:G1} and the identity $Z_{k}/Z_{N} = \prod_{i > k}Y_{i}$.
    When $t$ has nonnegative exponents, $t/Z_{N}$ has zero pole order at every level and $\varphi_{k}(t/Z_{N}) = P_{k}(t) - s_{k}$.
    The criterion for $R_{0}$ gives~(b), even for a coefficient of negative value.
    The sum defining $S_{0}$ is a ring because $J_{0}$ is an ideal of $R_{0}$.
    Its two summands satisfy Lemma~\ref{lem:MC}, giving~(c) by~(b).
    Finally, in~(d), the quotient $w = t/Z_{N}$ lies in $R_{0}$ with $v(w) < 2\kappa$.
    Since $\varphi_{k}(w) \le v(w) < 2\kappa \le T_{k}$ for $k \ge 1$, $w$ can pass only at zero.
    Thus $e(w) \ge \kappa\delta_{0}(w) \ge 0$ and $e(t) = \kappa + e(w)$.
    If $e(w) = 0$, the pole order vanishes and $w$ is a polynomial monomial with unit coefficient.
\end{proof}

\begin{corollary}
    \label{rem:membership}
    At each fixed exponent, membership in any of $R_{0}, J_{0}, S_{0}$ depends only on the coefficient value.
\end{corollary}

\begin{proof}
    Multiplication by a unit of $\widehat{\mathcal{O}}$ preserves each of the three sets in both directions.
\end{proof}

The obstruction argument will compare two translations of a monomial, one for each part of a partition.
To control products after these translations we need a slightly larger ring.

\begin{lemma}
    \label{lem:G3}
    Let $\mathcal{A} \sqcup \mathcal{B} = \{2, \ldots, N\}$ and let $\mathcal{I}$ be the ideal of $R_{0}$ generated by the sets $\mathcal{H}_{k}$.
    Set
    \[
        R_{0}^{\sharp} = R_{0} + \varpi^{-\kappa}\mathcal{I},
        \qquad U_{0} = \{f \in R_{0}^{\sharp} : Y_{\mathcal{B}}f \in R_{0}\}.
    \]
    \begin{enumerate}
        \item[\textup{(a)}] Both are rings containing $R_{0}$, with $Y_{\mathcal{B}}U_{0} \subseteq R_{0}$ and $\varpi^{-2\kappa}\mathcal{I}^{2} \subseteq \mathcal{I}$.
        \item[\textup{(b)}] A monomial $t$ belongs to $U_{0}$ whenever both $Y_{\mathcal{A}}t$ and $Y_{\mathcal{B}}t$ belong to $R_{0}$.
    \end{enumerate}
\end{lemma}

\begin{proof}
    For $h \in \mathcal{H}_{k}$ and $h' \in \mathcal{H}_{j}$ with $j \le k$, the monomial $\varpi^{-2\kappa}hh'$ involves only the first $k$ variables and has value at least $T_{k} + T_{j} - 2\kappa \ge T_{k}$.
    Thus $M = \varpi^{-\kappa}\mathcal{I}$ satisfies $M^{2} \subseteq \mathcal{I} \subseteq R_{0}$, and $R_{0}^{\sharp} = R_{0} + M$ is a ring.
    Since $Y_{\mathcal{B}} \in R_{0}$, one has $U_{0} = R_{0} + (M \cap Y_{\mathcal{B}}^{-1}R_{0})$.
    Here the second summand is an $R_{0}$-module whose square lies in $R_{0}$.
    This proves~(a).

    For (b), multiplying by $Y_{X}$, for $X = \mathcal{A}$ or $\mathcal{B}$, lowers any pole order by at most one.
    At $k \ge 1$ this yields
    \[
        \varphi_{k}(Y_{X}t) - \varphi_{k}(t)
        \le u_{X_{\le k}} + s_{k} \le s_{1} < \kappa,
    \]
    since $1 \notin X$.
    If either translate passes at a positive level $k$, then $\varpi^{\kappa}t$ passes there as well.
    Its factorization in Lemma~\ref{lem:G1}(b) contains a factor in $\mathcal{H}_{k}$, so $t \in M$.
    Otherwise both translates pass at level zero.
    A coordinate realizing the largest pole of $t$ is unchanged by at least one translate, since the two index sets are disjoint.
    Hence the larger of the two translated pole orders at zero equals $\delta_{0}(t)$, including when this is zero.
    As neither translation changes the coefficient, $e(t) \ge \kappa\delta_{0}(t)$, so $t \in R_{0}$.
    In both cases $t \in R_{0}^{\sharp}$, and $Y_{\mathcal{B}}t \in R_{0}$ completes the proof.
\end{proof}

\subsection{A product with prescribed initial term}

\begin{theorem}
    \label{thm:F}
    In the finite model, every factorization of $Z_{1}(1 + h)$ in $S_{0}$, where $h \in J_{0}$, has a factor in $J_{0}$.
\end{theorem}

\begin{proof}
    Suppose that $f_{1} \cdots f_{r} = Z_{1}(1 + h)$ with $r \ge 1$ and every $f_{i} \in S_{0} \setminus J_{0}$.
    Since $v(h) > 0$, the product has a unique least-value term whose exponent is that of $Y_{1}^{-1}$ and whose coefficient has value $\kappa$.
    By Lemma~\ref{lem:V} each factor has a unique least-value term $m_{i}$, and both the exponents and the coefficient values of these terms add.
    Their Gauss values are nonnegative and sum to $s_{1}$, so each is less than $2\kappa$.
    Some $m_{i}$ has a negative first exponent and hence belongs to $J_{0}$.
    By Lemma~\ref{lem:G2} each $m_{i}$ in $J_{0}$ has coefficient value at least $\kappa$, and each $m_{i}$ outside $J_{0}$ has nonnegative coefficient value.
    Thus precisely one $m_{i}$ belongs to $J_{0}$.
    Renumber so this is $m_{1}$.
    The same lemma writes $m_{1}$ as a unit coefficient times $Z_{N}Y^{\beta}$ with $\beta \ge 0$.
    Each other $m_{i}$ is a unit coefficient times an ordinary monomial supported on $\Lambda$.

    After removing $Z_{N}$, the exponents to be distributed among the $m_{i}$ are those of $Y_{2} \cdots Y_{N}$.
    Consequently there is a partition $\mathcal{A} \sqcup \mathcal{B} = \{2, \dots, N\}$ such that $Y^{\beta} = Y_{\mathcal{B}}$ and $f_{2} \cdots f_{r}$ has initial term $cY_{\mathcal{A}}$, with $v(c) = 0$ and $\mathcal{A} \subseteq \Lambda$.
    The empty product is $1$.
    In particular $2u_{\mathcal{A}} \le s_{1}$.
    Define $d$ by $f_{2} \cdots f_{r} = cY_{\mathcal{A}}(1 - d)$.
    Then $v(d) > 0$ and $Y_{\mathcal{A}}d \in S_{0}$.
    In $cf_{1}$ collect the terms outside $J_{0}$ into $a$, and write
    \[
        cf_{1} = a + Z_{N}\rho,\qquad \rho \in R_{0},\qquad \psi = a/Z_{N}.
    \]
    The termwise membership tests give $a \ne 0$ and put every term of $\psi$ outside $R_{0}$.
    Since $Z_{1} = Z_{N}Y_{\mathcal{A}}Y_{\mathcal{B}}$, the proposed factorization gives
    \[
        \psi + \rho = Y_{\mathcal{B}}\frac{1 + h}{1 - d}.
    \]

    We claim that a term of least value in $\psi$ must be matched, at the same exponent and value, by a term of $Y_{\mathcal{B}}d$.
    To see this, let $g$ be the sum of the terms of $d$ in the auxiliary ring $U_{0}$ of Lemma~\ref{lem:G3}, and put $b = d - g$.
    Every term $t$ of $b$ satisfies $Y_{\mathcal{A}}t \in S_{0}$ by termwise membership.
    Since $t \notin U_{0}$, Lemma~\ref{lem:G3} puts $Y_{\mathcal{B}}t$ outside $R_{0}$.
    Set $\nu = \min\{v(\psi), u_{\mathcal{B}} + v(b)\}$, which is finite.
    In the fraction field of the Laurent polynomial ring we have the exact identity
    \[
        \frac{1}{1 - d} - \frac{1}{1 - g} - b
        = \frac{b(d + g - dg)}{(1 - d)(1 - g)}.
    \]
    Both denominators have value zero and $d + g - dg$ has positive value.
    The geometric sum $G_{n} = \sum_{j = 0}^{n - 1}g^{j}$ belongs to $U_{0}$, and $(1 - g)^{-1} - G_{n} = g^{n}/(1 - g)$ has arbitrarily large value as $n$ increases.
    Hence for sufficiently large $n$,
    \[
        v\bigl(\psi - Y_{\mathcal{B}}b - H_{n}\bigr) > \nu,
        \qquad H_{n} = Y_{\mathcal{B}}(1 + h)G_{n} - \rho \in R_{0}.
    \]
    Indeed the displayed difference is the sum of $Y_{\mathcal{B}}hb$, the geometric truncation error times $Y_{\mathcal{B}}(1 + h)$, and the right-hand side of the preceding identity times $Y_{\mathcal{B}}(1 + h)$.
    The estimate includes $b = 0$ or $g = 0$ under the convention $v(0) = +\infty$.

    All terms of $\psi$ and $Y_{\mathcal{B}}b$ lie outside $R_{0}$.
    If a term of value $\nu$ in either polynomial had no match of that value at its exponent in the other, their difference would retain the same exponent and coefficient value.
    Since the difference agrees with $H_{n}$ up to terms of value greater than $\nu$, termwise membership would put that term in $R_{0}$.
    Thus the least terms match, both polynomials have value $\nu$, and $b \ne 0$.

    Choose a least-value term $\ell$ of $a$.
    Its translate $\ell/Z_{N}$ matches a term $Y_{\mathcal{B}}t$ with $t$ a term of $b$.
    Since $Y_{\mathcal{A}}t \in S_{0}$ and membership depends only on exponent and coefficient value, it follows that
    \[
        \sigma = \frac{\ell Y_{\mathcal{A}}}{Z_{N}Y_{\mathcal{B}}}
        = \frac{\ell Y_{\mathcal{A}}^{2}}{Z_{1}} \in S_{0}.
    \]
    But $\ell \notin J_{0}$ is an ordinary monomial with $P_{k}(\ell) < \tau_{k}$ at every level.
    In particular $e(\sigma) = e(\ell) - \kappa < 0$, excluding $\sigma$ from the polynomial part of $S_{0}$.
    Its exponents are nonnegative, so Lemma~\ref{lem:G2} requires $P_{k}(\sigma) \ge \tau_{k}$ at some level for $\sigma$ to lie in $J_{0}$.
    Level zero is excluded by $e(\sigma) < 0$.
    For every $k \ge 1$,
    \[
        P_{k}(\sigma) = P_{k}(\ell) - s_{1} + 2u_{\mathcal{A}_{\le k}}
        \le P_{k}(\ell) < \tau_{k}.
    \]
    This contradiction proves the theorem.
\end{proof}

\section{Evaluation on the algebraic blocks}
\label{sec:transfer}

Take the blocks of Proposition~\ref{prop:B} and fix a set $\Lambda$ of positive integers for which $2\sum_{i \in \Lambda,\, i > 1} u_{i} \le s_{1}$.
At stage $N$ use the finite model for the first $N$ blocks, selecting the indices in $\Lambda \cap [1, N]$.
Evaluation at $(y_{1}, \dots, y_{N})$ sends $Z_{i}$ to $z_{i}$ and preserves the value of each monomial.
It can raise the value of a sum, since the $y_{i}$ satisfy algebraic relations.
We construct an algebra of representatives on which every such relation belongs to $Z_{1}J_{0}$.

\begin{definition}
    \label{def:R}
    For $l \ge 1$, $e \in \mathbb{Z}_{\ge 0}$ and $\mu \in \mathbb{Z}_{\ge 0}^{l}$, consider the inequalities
    \[
        e + \sum_{j = 1}^{i} \mu_{j}u_{j} \ge \tau_{i}
        \qquad (0 \le i \le l).
    \]
    If $l$ is the first index for which the inequality holds, adjoin the monomial
    \[
        Q_{e, \mu, l} = \frac{\varpi^{e}Y^{\mu}}{Z_{l}}
        = \varpi^{e - \kappa}\prod_{j = 1}^{l}Y_{j}^{\mu_{j} + 1}.
    \]
    Write $\rho_{e, \mu, l}$ for its evaluation.
    In $k_{0}[Y_{i}^{\pm 1} : i \ge 1]$, let $\mathcal{E}_{N}$ be the $\mathcal{O}$-algebra on the $Y_{i}, Z_{i}$ with $i \le N$ and all the $Q_{e, \mu, l}$ with $l \le N$.
    Set $\mathcal{E} = \bigcup_{N \ge 1}\mathcal{E}_{N}$ and $\mathcal{K}_{N} = \ker(\operatorname{ev}|_{\mathcal{E}_{N}})$, where $\operatorname{ev}$ sends each $Y_{i}$ to $y_{i}$ and fixes $k_{0}$.
    In $K$ define
    \[
        R_{*} = \operatorname{ev}(\mathcal{E}), \qquad
        J_{*} = \bigcup_{l \ge 1} z_{l}R_{*}, \qquad
        S_{*}^{\Lambda} = \mathcal{O}[y_{i} : i \in \Lambda] + J_{*}.
    \]
\end{definition}

\begin{lemma}
    \label{lem:Rclaim}
    The ring $R_{*}$ lies in $A \cap V$, and $J_{*}$ is an ideal of $R_{*}$.
    In particular $S_{*}^{\Lambda}$ is a subring of $A \cap V$ containing $\mathcal{O}$.
\end{lemma}

\begin{proof}
    For a generator $Q = Q_{e, \mu, l}$ the inequality at level $l$ gives
    \[
        v(Q) = v(\rho_{e, \mu, l}) = e + \sum_{j = 1}^{l}\mu_{j}u_{j} - s_{l}
        \ge \tau_{l} - s_{l} = T_{l}.
    \]
    Lemma~\ref{lem:D} clears the denominator $\varpi^{\kappa}$ in its evaluation, so $\rho_{e, \mu, l} \in A \cap V$.
    The other generators evaluate to $y_{i}, z_{i} \in A \cap V$.
    The identity $z_{l} = y_{l + 1}z_{l + 1}$ makes $J_{*}$ an increasing union of ideals.
\end{proof}

\begin{lemma}
    \label{lem:N}
    If $g \in \mathcal{O}[Y_{1}, \dots, Y_{N}] \cap J_{0}$, then $g(y_{1}, \dots, y_{N}) \in J_{*}$.
\end{lemma}

\begin{proof}
    By Lemma~\ref{lem:G2}, each term $cY^{\mu}$ of $g$ reaches the threshold $\tau_{i}$ at some index $i$.
    Put $e = v(c)$ and take the first such $i$.
    For $i = 0$ the coefficient belongs to $\varpi^{\kappa}\mathcal{O}$, and $\varpi^{\kappa} = y_{1}z_{1} \in J_{*}$.
    For $i > 0$ the choice of generators gives
    \[
        cy^{\mu}
        = (c\varpi^{-e})\left(\prod_{j = i + 1}^{N} y_{j}^{\mu_{j}}\right)
          z_{i}\rho_{e, \mu_{\le i}, i} \in z_{i}R_{*}.
    \]
    Thus all the evaluated terms belong to $J_{*}$.
\end{proof}

\begin{lemma}
    \label{lem:X}
    At every finite stage $N$, we have $\mathcal{E}_{N} \subseteq R_{0}$ and $\mathcal{K}_{N} \subseteq Z_{1}J_{0}$.
\end{lemma}

\begin{proof}
    Each $Z_{i}$ with $i \le N$ is $Z_{N}\prod_{i < j \le N}Y_{j}$, and each $Q_{e, \mu, i}$ uses only the first $i$ variables and has value at least $T_{i}$.
    The defining generators of $\mathcal{E}_{N}$ therefore belong to $R_{0}$.

    For the assertion about the kernel, first consider one product of these generators, written
    \[
        t = a\prod_{i = 1}^{N}Y_{i}^{p_{i}} Z_{i}^{b_{i}}\prod_{r}Q_{r}
          = cY^{\beta},
        \qquad a \in \mathcal{O} \setminus \{0\},\quad p_{i}, b_{i} \ge 0.
    \]
    Here $Q_{r} = Q_{e_{r}, \mu^{r}, l_{r}}$ and both products have finite length.
    Suppose $\beta \notin B_{N}$ and let $l$ be the largest index with $\beta_{l} + \delta_{l}(\beta) \notin I_{l}$.
    The estimate will use the level $k = \max(\{l\} \cup \{l_{r} : r\})$ so that it counts the full value of every $Q_{r}$.
    Put $\delta = \delta_{l}(\beta)$, $\theta = \delta_{k}(\beta)$, $b_{\ge j} = \sum_{i = j}^{N}b_{i}$, $b_{>j} = \sum_{i = j + 1}^{N}b_{i}$ and $M_{j} = \sum_{r:\,l_{r} \ge j}(\mu^{r}_{j} + 1)$.
    Then $\beta_{j} = p_{j} + M_{j} - b_{\ge j}$, so $\delta \le b_{>l}$ and $\theta \le \min(\delta, b_{>k})$.
    Since $P_{k}(Z_{i}) = s_{\min(i, k)}$, the contribution of all the $Z_{i}$ to $P_{k}(t)$, less the pole correction $\theta s_{k}$, is
    \[
        \sum_{i = 1}^{k} b_{i}s_{i} + (b_{>k} - \theta)s_{k} \ge b_{\ge l} - \theta,
    \]
    using $s_{i} \ge 1$.
    The remaining factors give
    \begin{equation}
        \label{eq:transfer-cost}
        \varphi_{k}(t) \ge b_{\ge l} - \theta + p_{l}u_{l} + \sum_{r}v(Q_{r}).
    \end{equation}

    For a generator of level $l_{r} \ge l$, put $d_{r} = (\mu^{r}_{l} + 1)u_{l}$.
    Its value is at least $d_{r} - \kappa$ by its formula, and at least $T_{l_{r}} \ge 2\kappa$ by its defining inequality.
    Taking two thirds of the first bound and one third of the second gives $v(Q_{r}) \ge 2d_{r}/3$.
    When $k > l$, one such generator has level $k$.
    The first threshold reached by this generator is at $k$, so $e_r + \sum_{j = 1}^{l}\mu^{r}_{j}u_{j} < \tau_{l}$ and in particular $\mu^{r}_{l}u_{l} < \tau_{l}$.
    Consequently $d_{r} < \tau_{l} + u_{l} = T_{l} + s_{l - 1} \le 3T_{l}/2$, and this generator contributes at least $T_{k} \ge T_{k} - T_{l} + 2d_{r}/3$.
    The generators of level below $l$ have positive value.
    Summing these bounds, including the case $k = l$, yields
    \[
        \sum_{r}v(Q_{r}) \ge T_{k} - T_{l} + \frac{2}{3} M_{l}u_{l}.
    \]

    The failed window now gives a lower bound at either end of $I_{l}$.
    If $\beta_{l} + \delta < -L_{l}$, then $b_{\ge l} - \theta \ge b_{\ge l} - \delta \ge L_{l} + 1 > D_{l}/2$.
    If $\beta_{l} + \delta \ge D_{l} - L_{l}$, then $p_{l} + M_{l} = \beta_{l} + b_{\ge l} \ge D_{l}/2$, since $b_{\ge l} \ge \delta$.
    In the first case \eqref{eq:transfer-cost} gives at least $T_{k} - T_{l} + D_{l}/2$, and in the second at least $T_{k} - T_{l} + D_{l}u_{l}/3$.
    As $u_{l} < \kappa$ and $\kappa \ge 1$, Proposition~\ref{prop:B} gives the common bound
    \begin{equation}
        \label{eq:transfer-threshold}
        \varphi_{k}(t) \ge T_{k} - T_{l} + \frac{D_{l}u_{l}}{3\kappa}
        \ge T_{k} + T_{l} + \Omega_{l}
        \ge \tau_{k} + s_{1} + \Omega_{l}.
    \end{equation}
    The last inequality uses $s_{k} + s_{1} \le 2\kappa \le T_{l}$.

    Let $\sigma_{N}$ denote the $\widehat{k}_{0}$-linear map from $F_{N}$ to the Laurent polynomial ring sending the window basis $y^{\gamma}$ to $Y^{\gamma}$ for $\gamma \in B_{N}$.
    For the product $t$ above set
    \[
        m = Z_{l}^{\delta}\prod_{i = l + 1}^{N}Y_{i}^{\beta_{i}},
        \qquad h = t/m.
    \]
    The monomial $h$ uses only the first $l$ variables.
    Write $h(y) = \sum_{\gamma \in B_{l}} c_{\gamma}y^{\gamma}$ in the basis of Lemma~\ref{lem:WB}.
    For every nonzero coefficient, Lemma~\ref{lem:WL} gives $v(c_{\gamma}Y^{\gamma}) \ge v(h) - \Omega_{l}$.
    By Lemma~\ref{lem:W}, the exponents $(\gamma - \delta\mathbf{1},\,\beta_{>l})$ of $Y^{\gamma}m$ all belong to $B_{N}$, since the tail passes every window above $l$.
    They are distinct, so
    \[
        \sigma_{N}(t(y)) = \sum_{\gamma \in B_{l}} c_{\gamma}Y^{\gamma}m.
    \]
    Multiplication of $m$ by any monomial in the first $l$ variables adds that monomial's value to $\varphi_{k}(m)$, since $k \ge l$ and the pole order above $k$ stays fixed.
    Thus every term $x$ on the right, as well as $t$ itself, satisfies $\varphi_{k}(x) \ge \tau_{k} + s_{1}$ by \eqref{eq:transfer-threshold}.
    Dividing $x$ by $Z_{1}Z_{N}$ lowers $P_{k}(x)$ by $s_{1} + s_{k}$ and can only decrease the pole order above $k$.
    Hence $\varphi_{k}(x/(Z_{1}Z_{N})) \ge T_{k}$, so Lemma~\ref{lem:G1} gives $x \in Z_{1}Z_{N}R_{0} = Z_{1}J_{0}$.

    We have proved $t - \sigma_{N}(t(y)) \in Z_{1}J_{0}$ for every product of generators whose exponent fails a window.
    For exponents in $B_{N}$ this difference is zero by definition.
    Every $\Pi \in \mathcal{K}_{N}$ is a sum of such products and satisfies $\sigma_{N}(\Pi(y)) = 0$.
    Addition gives $\Pi \in Z_{1}J_{0}$.
\end{proof}

\section{The element without an atomic factorization}
\label{sec:witness}

\begin{proof}[Proof of Theorem~\ref{thm:dvr}]
    Choose a strictly ascending chain of nonzero principal ideals in $R$ and apply Lemma~\ref{lem:normalize} with $A = R$.
    For the blocks of Proposition~\ref{prop:B}, the partial sums of $\sum_{i \ge 1}u_i$ equal $\kappa - s_j$ at stage $j$ and are bounded by $\kappa - 1$.
    Choose an integer $m \ge 1$ with $2\sum_{i > m}u_i \le s_1$ and take $\Lambda = \{i : i > m\}$ in Definition~\ref{def:R}.
    This gives $S = S_{*}^{\Lambda}$ with $\mathcal{O} \subseteq S \subseteq R \cap V$.
    The element $z_1 \in J_{*}$ is a nonzero nonunit of $S$ since $v(z_1) = s_1 \ge 1$.

    For nonzero $f \in J_{*}$, write $f = z_j g$ with $g \in R_{*}$.
    Pick $l \in \Lambda$ with $l > j$.
    The identities between successive blocks give a factorization in $S$,
    \[
        f = y_l \left(z_l g\prod_{j < i < l} y_i\right).
    \]
    The second factor belongs to $J_{*}$.
    Both factors have positive value, so neither is a unit and $f$ is reducible.
    Every atom of $S$ therefore lies outside $J_{*}$.

    Suppose $z_1 = f_1 \cdots f_r$ with $r \ge 1$ and each $f_i \in S \setminus J_{*}$.
    Write $f_i = p_i(y) + z_{l_i} g_i$, where $p_i$ is a polynomial over $\mathcal{O}$ in the selected variables and $g_i \in R_{*}$.
    Each $g_i$ is represented by an element $G_i \in \mathcal{E}$.
    There is a finite stage $N$ containing all the indices used to form the $p_i$, $Z_{l_i}$, and $G_i$.
    At this stage put $\Phi_i = p_i(Y) + Z_{l_i} G_i$.
    Lemma~\ref{lem:X} puts $G_i$ in $R_{0}$, so $Z_{l_i} G_i \in J_{0}$ and $\Phi_i \in S_{0}$.
    Since $f_i \notin J_{*}$, also $p_i(y) \notin J_{*}$.
    Lemma~\ref{lem:N} gives $p_i(Y) \notin J_{0}$, hence $\Phi_i \notin J_{0}$.

    Now $\Phi_1 \cdots \Phi_r - Z_1$ belongs to $\mathcal{E}_{N}$ and evaluates to zero.
    The kernel inclusion of Lemma~\ref{lem:X} yields an identity
    \[
        \Phi_1 \cdots \Phi_r = Z_1(1 + h), \qquad h \in J_{0}.
    \]
    Theorem~\ref{thm:F} forces one of the $\Phi_i$ into $J_{0}$, a contradiction.
    Every factorization of $z_1$ in $S$ thus has a factor in $J_{*}$, and every such nonzero factor can be split again.
    In particular $z_1$ has no factorization into atoms.
\end{proof}

\begin{remark}
    \label{rem:nonprime}
    The ideal $J_{*}$ is not prime in $S$.
    Indeed $s_j = s_{j + 1} + u_{j + 1} > 1$ for every $j$, so every nonzero element of $J_{*}$ has value strictly greater than one.
    Thus $\varpi \notin J_{*}$, whereas $\varpi^{\kappa} = y_1 z_1 \in z_1 R_{*} \subseteq J_{*}$.
    The product of $\kappa$ copies of $\varpi$ shows that the property proved for $z_1$ fails for another element of the same ideal.
\end{remark}

\section{A witness in the semilocal domain}
\label{sec:semilocal}

Let $D$ be a semilocal domain and write $K = \Frac(D)$.
A chain in one localization can involve quotients that lie outside $D$.
We first replace its generators by elements of $D$ that are units at all the other maximal ideals.
This makes it possible to split elements of the constructed ideal after intersecting the local subring with $D$.

\begin{proposition}
    \label{prop:semilocal-witness}
    Let $M$ be a maximal ideal of $D$ and suppose that $D_{M}$ contains a discrete valuation ring $\mathcal{O}$ with a uniformizer $\varpi \in D$.
    Suppose also that $K$ is algebraic over $\Frac(\mathcal{O})$.
    If $D_{M}$ fails ACCP, then $D$ contains a nonatomic subring.
\end{proposition}

\begin{proof}
    Put $R = D_{M}$ and $\mathfrak{m} = MR$.
    Since $R$ fails ACCP, it is a nonfield, and the proof of Lemma~\ref{lem:normalize} shows that $\mathfrak{m}$ is its only nonzero prime.
    We first lift each nonzero $\xi \in R$ to an element $\eta \in D$ that generates $\xi R$ in $R$ and is a unit in $D_{P}$ at every maximal ideal $P \ne M$.
    Write $\xi = a/b$ with $a \in D$ and $b \in D \setminus M$.
    The ideal $a\mathfrak{m}$ of $R$ is nonzero and proper, so $\sqrt{a\mathfrak{m}} = \mathfrak{m}$.
    Its contraction $I = a\mathfrak{m} \cap D$ has radical $M$ and is therefore comaximal with every other maximal ideal of $D$.
    By the Chinese remainder theorem choose $\eta \in D$ congruent to $a$ modulo $I$ and to $1$ modulo every maximal ideal $P \ne M$.
    Then $\eta/a \in 1 + \mathfrak{m}$, so $\eta R = \xi R$, as required.

    Lift generators of a strictly ascending chain of nonzero principal ideals of $R$ in this way, obtaining $\eta_{0}, \eta_{1}, \dots$ in $D$.
    The quotients $t_{i} = \eta_{i - 1}/\eta_{i}$ belong to $\mathfrak{m}$ and are units in every $D_{P}$ with $P \ne M$.
    Thus $t_{i} \in D$, since $D$ is the intersection in $K$ of its localizations at maximal ideals.
    Indeed, if $x$ belongs to every such localization, its denominator ideal $\{d \in D : dx \in D\}$ meets the complement of every maximal ideal and hence is $D$.

    Apply Lemma~\ref{lem:normalize} and the construction in the proof of Theorem~\ref{thm:dvr} to the chain $\eta_{i}R$.
    With the notation of Proposition~\ref{prop:B}, all the elements
    \[
        y_{j} = \prod_{a_{j - 1} < i \le a_{j}} t_{i},
        \qquad z_{j} = \frac{\varpi^{\kappa}}{y_{1}\cdots y_{j}}
    \]
    belong to $R$.
    At each other maximal ideal $P$ the $y_{j}$ are units of $D_{P}$, and $\varpi \in D$ gives $z_{j} \in D_{P}$ as well.
    Consequently $y_{j}, z_{j} \in D$ for all $j$.

    Take $S = S_{*}^{\Lambda}$ and $J_{*}$ from that construction and put $S' = S \cap D$.
    The local proof gives $S \subseteq V$ and a positive-valued element $z_{1} \in S'$ for which every factorization in $S$ has a factor in $J_{*}$.
    Consider a nonzero $f \in J_{*} \cap D$.
    Write $f = z_{j}g$ with $g \in R_{*}$, and choose $l \in \Lambda$ with $l > j$.
    Then $y_{l} \in S'$ and
    \[
        \frac{f}{y_{l}} = z_{l}g\prod_{j < i < l} y_{i} \in J_{*} \subseteq R.
    \]
    At all the other maximal ideals, division by the unit $y_{l}$ preserves membership in $D_{P}$.
    Thus $f/y_{l}$ belongs to $J_{*} \cap D$.
    Both $y_{l}$ and $f/y_{l}$ have positive value and are nonunits of $S'$.
    Every factorization of $z_{1}$ in $S'$ is a factorization in $S$, hence includes a factor in $J_{*} \cap D$ which splits as above.
    This proves that $z_{1}$ has no factorization into atoms of $S'$.
\end{proof}

\begin{lemma}
    \label{lem:semilocal-base}
    Suppose that $D$ is hereditarily atomic and is not a field.
    In characteristic zero, $K$ is algebraic over $\mathbb{Q}$.
    In characteristic $p > 0$, there is an element $\pi \in D$ transcendental over $\mathbb{F}_{p}$ for which $K$ is algebraic over $\mathbb{F}_{p}(\pi)$.
    For every maximal ideal $M$ of $D$, there is a discrete valuation ring $\mathcal{O}_{M} \subseteq D_{M}$ with fraction field equal to the respective base field and with a uniformizer in $D$.
\end{lemma}

\begin{proof}
    Suppose first that $D$ has characteristic zero.
    Only finitely many rational primes are nonunits of $D$, since the contraction of each maximal ideal to $\mathbb{Z}$ contains at most one rational prime.
    Choose a rational prime $q$ outside all of them.
    It is a unit of $D$ and a nonunit of $\mathbb{Z}$.
    Lemma~\ref{lem:transcendental-witness} and hereditary atomicity show that every element of $D$ is algebraic over $\mathbb{Q}$, so the same holds for $K$.
    If $M \cap \mathbb{Z} = 0$, then $D_{M}$ contains $\mathbb{Q}$ and is algebraic over it.
    Such a domain is a field, whereas $MD_{M}$ is nonzero since $D$ is a nonfield domain.
    Thus $M \cap \mathbb{Z} = p\mathbb{Z}$ for a rational prime $p$, and $\mathcal{O}_{M} = \mathbb{Z}_{(p)} \subseteq D_{M}$ has uniformizer $p \in D$.

    Now let $D$ have characteristic $p > 0$.
    Its Jacobson radical contains a nonzero element $\pi$, obtained by multiplying a nonzero element from each maximal ideal.
    If $\pi$ were algebraic over $\mathbb{F}_{p}$, the ring $\mathbb{F}_{p}[\pi]$ would be a field and $\pi$ a unit of $D$.
    Thus $A = \mathbb{F}_{p}[\pi]$ is a polynomial ring and $1 + \pi$ is a nonunit of $A$.
    It is a unit of $D$, so Lemma~\ref{lem:transcendental-witness} gives algebraicity over $\mathbb{F}_{p}(\pi)$ as before.
    Every polynomial in $\pi$ with nonzero constant coefficient maps to a nonzero scalar modulo every maximal ideal of $D$.
    Hence $A_{(\pi)} \subseteq D$ is a single DVR that works for all $M$.
\end{proof}

\begin{proposition}
    \label{prop:semilocal-accp}
    If $D$ is a semilocal hereditarily atomic domain, then $D$ and every localization $D_{M}$ at a maximal ideal satisfy ACCP.
\end{proposition}

\begin{proof}
    The assertion for fields is immediate.
    Otherwise Lemma~\ref{lem:semilocal-base} supplies a base for Proposition~\ref{prop:semilocal-witness} at each maximal ideal $M$.
    Since every subring of $D$ is atomic, each $D_{M}$ satisfies ACCP.
    If $D$ admitted a strictly ascending chain of nonzero principal ideals, the ratios of successive generators would be nonunits.
    Finiteness of the set of maximal ideals gives a maximal ideal $M$ containing infinitely many of those ratios.
    At each of these indices the localized chain would have a strict inclusion, contradicting ACCP in $D_{M}$.
\end{proof}

\section{Valuations that detect units of the semilocal domain}
\label{sec:places}

Let $D$ be a semilocal hereditarily atomic domain that is not a field, and put $K = \Frac(D)$.
By Lemma~\ref{lem:semilocal-base}, $K$ is algebraic over a global field $F$.
A unit of $D$ can be a nonunit of a subring $S$.
To detect such an element one can dominate a localization of $S$ in which it is a nonunit.
The resulting valuation ring omits its inverse and therefore omits part of $D$.
We determine how many such valuations can occur on divisors of a fixed element.

We take $F = \mathbb{Q}$ or $F = \mathbb{F}_p(\pi)$ as in that lemma.
Every nontrivial valuation of $K$ has rank one, meaning that its value group embeds as an ordered subgroup of $\mathbb{R}$.
To see this, let $v$ be such a valuation.
For $x \in K^{\times}$, two nonzero terms in its minimal polynomial relation over $F$ have equal least value.
Consequently some positive multiple of $v(x)$ belongs to $v(F^{\times})$.
The restriction to $F$ is nontrivial, and its value group is cyclic.
Indeed, a nontrivial valuation on $\mathbb{Q}$ is nonnegative on $\mathbb{Z}$ and positive on exactly one rational prime.
In the rational function case it is trivial on $\mathbb{F}_p^{\times}$.
If $v(\pi) < 0$, the leading term gives the value of each nonzero polynomial in $\pi$.
If $v(\pi) \ge 0$, the contraction of the maximal ideal to $\mathbb{F}_p[\pi]$ is a nonzero prime, and every polynomial relatively prime to its generator has value zero.
In each case, normalizing the positive generator of $v(F^{\times})$ to one identifies $v(K^{\times})$ with an ordered subgroup of $\mathbb{Q}$.

\begin{lemma}
    \label{lem:simultaneous-units}
    Let $v_1, \dots, v_s$ be pairwise inequivalent nontrivial valuations of $K$ whose valuation rings do not contain $D$.
    Given $\gamma_i \in v_i(K^{\times})$ for each $i$, there is a unit $u \in D^{\times}$ with $v_i(u) = \gamma_i$ for every $i$.
\end{lemma}

\begin{proof}
    By weak approximation there is $x \in K^{\times}$ with $v_i(x) = \gamma_i$ for every $i$.
    Indeed, choose $x_i$ of value $\gamma_i$ at $v_i$ and approximate so that $v_i(x - x_i) > \gamma_i$.
    The Artin--Whaples approximation theorem applies to the inequivalent absolute values $x \mapsto \exp(-v_i(x))$ \cite[Theorem~1.1.3]{MR2183496}.

    To obtain all these values on a unit, choose, for each $i$, an element $d_i \in D$ with $v_i(d_i) < 0$.
    There is a positive integer $t$ such that at each valuation $v_i$ the negative numbers among $t^j v_i(d_j)$, $1 \le j \le s$, are pairwise distinct.
    Each possible tie between two negative entries excludes at most one positive $t$.
    Then $d = \sum_{j = 1}^{s}d_j^{t^j}$ has negative value at every $v_i$, since the least value occurs only once.
    Choose nonzero $a, b \in D$ with $x = a/b$ and nonzero $h \in \bigcap_{M \in \operatorname{Max}(D)}M$.
    Such $h$ exists by taking a product of one nonzero element from each maximal ideal.
    Since the valuations have rank one, there is an integer $n$ large enough that both $hd^n a$ and $hd^n b$ have negative value at every $v_i$.
    Their sums with $1$ are units of $D$, so
    \[
        u = \frac{1 + hd^n a}{1 + hd^n b} \in D^{\times},
        \qquad v_i(u) = v_i(a) - v_i(b) = \gamma_i. \qedhere
    \]
\end{proof}

\begin{proposition}
    \label{prop:outside}
    Every valuation ring $W$ of $K$ which does not contain $D$ is a DVR.
    Each $z \in K^{\times}$ is a unit of all but finitely many such $W$.
\end{proposition}

\begin{proof}
    Let $v$ have valuation ring $W$ with $D \nsubseteq W$.
    Lemma~\ref{lem:simultaneous-units} gives $v(D^{\times}) = v(K^{\times})$.
    In the atomic subring $T = D \cap W$, factor a $D$-unit of any prescribed positive value $\gamma$ into atoms.
    Each factor $a$ is a $D$-unit with $v(a) > 0$, since a factor of value zero has its inverse in $T$.
    If the value group contained $0 < \beta < v(a)$, a $D$-unit $b$ of value $\beta$ would split $a = b(a/b)$ into two positive-valued elements of $T$.
    Thus $v(a)$ is the least positive value, and every positive $\gamma$ is a sum of copies of it.
    The value group is discrete cyclic, so $W$ is a DVR.

    For the second assertion, suppose that some $z \in K^{\times}$ has nonzero value at infinitely many distinct valuation rings omitting part of $D$.
    Pass to a sequence of such rings on which the sign of the value of $z$ is constant.
    The field $K$ is countable, since it is algebraic over the countable field $F$.
    Enumerate $K^{\times}$ and successively pass to infinite subsequences on which the sign of the value of the next element is constant.
    Taking a diagonal subsequence gives pairwise distinct rings along which the sign for every fixed element is eventually constant.

    The elements with eventually nonnegative value, together with zero, form a subring $W_{\infty}$ by the valuation inequalities.
    For each nonzero $x$, at least one of $x, x^{-1}$ belongs to it, so $W_{\infty}$ is a valuation ring.
    An element has positive value there precisely when its values along the sequence are eventually positive, and it is a unit there precisely when they are eventually zero.
    The sign chosen for $z$ makes $W_{\infty}$ nontrivial.
    Discard $W_{\infty}$ itself if it occurs in the sequence.

    Choose nonzero $h \in D$ of positive value at $W_{\infty}$.
    If $D \subseteq W_{\infty}$, some such $h$ exists since otherwise every nonzero element of $D$ and consequently of $K$ would be a unit there.
    If $D \nsubseteq W_{\infty}$, Lemma~\ref{lem:simultaneous-units} even gives such $h$ in $D^{\times}$.
    After discarding finitely many terms, $h$ has positive value throughout the sequence.

    We select a subsequence $W_1, W_2, \dots$ with normalized valuations $w_i$ and construct units $u_n \in D^{\times}$ such that
    \[
        w_n(u_n) = 1, \qquad w_i(u_n) = 0 \ (i \ne n),
        \qquad u_n \in W_{\infty}^{\times}.
    \]
    Having chosen the first $n - 1$ pairs, choose $W_n$ sufficiently far along the sequence that each earlier $u_i$ is a unit there.
    Lemma~\ref{lem:simultaneous-units} prescribes the value one at $W_n$ and zero at the earlier chosen rings.
    If $W_{\infty}$ omits some element of $D$, include it in the prescription with value zero.
    If it contains $D$, it contains each $D$-unit and its inverse.
    This produces $u_n$, and each subsequent choice of ring preserves the required zero values.

    Consider the subring and its prime ideal
    \[
        S = D \cap W_{\infty} \cap \bigcap_{n \ge 1}W_n,
        \qquad P = S \cap \mathfrak{m}_{W_{\infty}}.
    \]
    We have $h \in P$.
    For any nonzero $y \in P$, all sufficiently large $n$ satisfy $w_n(y) > 0$, hence $w_n(y) \ge 1$.
    For such $n$ we have $u_n \in S$ and $y/u_n \in P$.
    Indeed, division by $u_n$ preserves membership in $D$, preserves the value at $W_{\infty}$ and all $W_i$ for $i \ne n$, and leaves a nonnegative value at $W_n$.
    Both factors of $y = u_n(y/u_n)$ are nonunits of $S$.
    Thus $P$ has no atoms, whereas a factorization of $h$ into atoms would place one of those atoms in $P$ by primality.
    This contradicts hereditary atomicity of $D$ and proves the finiteness assertion.
\end{proof}

\section{Passing ACCP to every subring}
\label{sec:hereditary}

\begin{lemma}
    \label{lem:unit-bound}
    Let $D$ be a nonfield semilocal hereditarily atomic domain and let $0 \ne z \in D$.
    There is an integer $B_D(z)$ such that, for every subring $S \subseteq D$ and every decomposition $z = b a_1 \cdots a_r$ with $b \in S$ and $a_i \in S \setminus S^{\times}$, at most $B_D(z)$ of the $a_i$ are units of $D$.
\end{lemma}

\begin{proof}
    Let $\mathcal{W}(z)$ be the valuation rings $W$ of $\Frac(D)$ such that $D \nsubseteq W$ and $z \in \mathfrak{m}_W$.
    By Proposition~\ref{prop:outside} this set is finite and each $W$ in it has a valuation $w_W$ normalized to have value group $\mathbb{Z}$.
    We claim that one can take
    \[
        B_D(z) = \sum_{W \in \mathcal{W}(z)} w_W(z).
    \]

    For each $a_i$ that is a unit of $D$, choose a maximal ideal $P$ of $S$ containing $a_i$.
    There is a valuation ring $W$ of $\Frac(D)$ dominating $S_P$ \cite{Stacks}.
    In particular $W$ contains $S$ and $a_i \in \mathfrak{m}_W$.
    Its inverse lies in $D$ and lies outside $W$, so $D \nsubseteq W$.
    All factors of $z$, including $b$, belong to $W$.
    Hence $z \in \mathfrak{m}_W$, and we can assign this index $i$ to $W \in \mathcal{W}(z)$.
    At any $W$ receiving an index every factor has nonnegative value, and each assigned factor has value at least one.
    Thus at most $w_W(z)$ indices can be assigned to $W$.
    Summing over $\mathcal{W}(z)$ gives the bound.
\end{proof}

\begin{proof}[Proof of Theorem~\ref{thm:hereditary}]
    A domain satisfying ACCP is atomic by the splitting argument in Section~\ref{sec:intro}.
    Applying this to each subring proves that hereditary ACCP implies hereditary atomicity.

    Suppose now that $D$ is hereditarily atomic.
    Proposition~\ref{prop:semilocal-accp} gives ACCP both for $D$ and for each $D_M$.
    Assume first that $D$ is not a field, and let $S$ be any subring of $D$.
    If $S$ fails ACCP, there are nonzero $x_0, x_1, \ldots \in S$ with $x_{i - 1} = a_i x_i$ and $a_i \in S \setminus S^{\times}$ for $i \ge 1$.
    The chain $x_0D \subseteq x_1D \subseteq \cdots$ stabilizes, so $a_i \in D^{\times}$ from some index onward.
    But for every $r$ we have $x_0 = x_r a_1 \cdots a_r$ in $S$.
    Lemma~\ref{lem:unit-bound} bounds the number of $D$-units among these $r$ quotients by the same integer $B_D(x_0)$, a contradiction.

    Finally suppose that $D$ is a field.
    A subring $S$ that is a field already satisfies ACCP.
    Otherwise choose a nonzero maximal ideal $P$ of $S$ and a valuation ring $V$ of $\Frac(S)$ dominating $S_P$ \cite{Stacks}.
    We have $S \subseteq V \subseteq \Frac(S) \subseteq D$.
    The maximal ideal of $V$ contains $P$, so $V$ is a nonfield local domain.
    Every subring of $V$ is atomic because it is a subring of $D$.
    The nonfield case just proved applies to $V$ and gives ACCP for its subring $S$.
\end{proof}

\begin{samepage}
    \paragraph{Authorship.}
    Astra 6.2 authored this manuscript in full.
    Claude Opus 5.5 discovered the DVR result (Theorem~\ref{thm:dvr}).
    The proofs received independent AI review.
    Human verification of the full proof is pending.
    \par
\end{samepage}

\begin{thebibliography}{9}
    \bibitem{CGH}
    J.~Coykendall, F.~Gotti and R.~E.~Hasenauer, Hereditary atomicity in integral domains,
    \emph{J. Pure Appl. Algebra} \textbf{227} (2023), no.~4, 107249.
    \href{https://doi.org/10.1016/j.jpaa.2022.107249}{doi:10.1016/j.jpaa.2022.107249}.

    \bibitem{MR2183496}
    A.~J.~Engler and A.~Prestel, \emph{Valued Fields},
    Springer Monographs in Mathematics, Springer, Berlin, 2005.
    \href{https://doi.org/10.1007/3-540-30035-X}{doi:10.1007/3-540-30035-X}.

    \bibitem{Grams}
    A.~Grams, Atomic rings and the ascending chain condition for principal ideals,
    \emph{Proc. Cambridge Philos. Soc.} \textbf{75} (1974), 321--329.
    \href{https://doi.org/10.1017/S0305004100048532}{doi:10.1017/S0305004100048532}.

    \bibitem{Stacks}
    The Stacks Project Authors, \emph{The Stacks Project}, Lemma~10.50.2.
    \href{https://stacks.math.columbia.edu/tag/00IA}{Tag 00IA}.
\end{thebibliography}

\end{document}
